% Options for packages loaded elsewhere
\PassOptionsToPackage{unicode}{hyperref}
\PassOptionsToPackage{hyphens}{url}
\PassOptionsToPackage{dvipsnames,svgnames,x11names}{xcolor}
%
\documentclass[
  12pt,
  a4paper,
  DIV=11,
  numbers=noendperiod]{scrartcl}

\usepackage{amsmath,amssymb}
\usepackage{iftex}
\ifPDFTeX
  \usepackage[T1]{fontenc}
  \usepackage[utf8]{inputenc}
  \usepackage{textcomp} % provide euro and other symbols
\else % if luatex or xetex
  \usepackage{unicode-math}
  \defaultfontfeatures{Scale=MatchLowercase}
  \defaultfontfeatures[\rmfamily]{Ligatures=TeX,Scale=1}
\fi
\usepackage[]{times}
\ifPDFTeX\else  
    % xetex/luatex font selection
\fi
% Use upquote if available, for straight quotes in verbatim environments
\IfFileExists{upquote.sty}{\usepackage{upquote}}{}
\IfFileExists{microtype.sty}{% use microtype if available
  \usepackage[]{microtype}
  \UseMicrotypeSet[protrusion]{basicmath} % disable protrusion for tt fonts
}{}
\makeatletter
\@ifundefined{KOMAClassName}{% if non-KOMA class
  \IfFileExists{parskip.sty}{%
    \usepackage{parskip}
  }{% else
    \setlength{\parindent}{0pt}
    \setlength{\parskip}{6pt plus 2pt minus 1pt}}
}{% if KOMA class
  \KOMAoptions{parskip=half}}
\makeatother
\usepackage{xcolor}
\usepackage[lmargin=1in,rmargin=1in]{geometry}
\setlength{\emergencystretch}{3em} % prevent overfull lines
\setcounter{secnumdepth}{5}
% Make \paragraph and \subparagraph free-standing
\makeatletter
\ifx\paragraph\undefined\else
  \let\oldparagraph\paragraph
  \renewcommand{\paragraph}{
    \@ifstar
      \xxxParagraphStar
      \xxxParagraphNoStar
  }
  \newcommand{\xxxParagraphStar}[1]{\oldparagraph*{#1}\mbox{}}
  \newcommand{\xxxParagraphNoStar}[1]{\oldparagraph{#1}\mbox{}}
\fi
\ifx\subparagraph\undefined\else
  \let\oldsubparagraph\subparagraph
  \renewcommand{\subparagraph}{
    \@ifstar
      \xxxSubParagraphStar
      \xxxSubParagraphNoStar
  }
  \newcommand{\xxxSubParagraphStar}[1]{\oldsubparagraph*{#1}\mbox{}}
  \newcommand{\xxxSubParagraphNoStar}[1]{\oldsubparagraph{#1}\mbox{}}
\fi
\makeatother


\providecommand{\tightlist}{%
  \setlength{\itemsep}{0pt}\setlength{\parskip}{0pt}}\usepackage{longtable,booktabs,array}
\usepackage{calc} % for calculating minipage widths
% Correct order of tables after \paragraph or \subparagraph
\usepackage{etoolbox}
\makeatletter
\patchcmd\longtable{\par}{\if@noskipsec\mbox{}\fi\par}{}{}
\makeatother
% Allow footnotes in longtable head/foot
\IfFileExists{footnotehyper.sty}{\usepackage{footnotehyper}}{\usepackage{footnote}}
\makesavenoteenv{longtable}
\usepackage{graphicx}
\makeatletter
\def\maxwidth{\ifdim\Gin@nat@width>\linewidth\linewidth\else\Gin@nat@width\fi}
\def\maxheight{\ifdim\Gin@nat@height>\textheight\textheight\else\Gin@nat@height\fi}
\makeatother
% Scale images if necessary, so that they will not overflow the page
% margins by default, and it is still possible to overwrite the defaults
% using explicit options in \includegraphics[width, height, ...]{}
\setkeys{Gin}{width=\maxwidth,height=\maxheight,keepaspectratio}
% Set default figure placement to htbp
\makeatletter
\def\fps@figure{htbp}
\makeatother
% definitions for citeproc citations
\NewDocumentCommand\citeproctext{}{}
\NewDocumentCommand\citeproc{mm}{%
  \begingroup\def\citeproctext{#2}\cite{#1}\endgroup}
\makeatletter
 % allow citations to break across lines
 \let\@cite@ofmt\@firstofone
 % avoid brackets around text for \cite:
 \def\@biblabel#1{}
 \def\@cite#1#2{{#1\if@tempswa , #2\fi}}
\makeatother
\newlength{\cslhangindent}
\setlength{\cslhangindent}{1.5em}
\newlength{\csllabelwidth}
\setlength{\csllabelwidth}{3em}
\newenvironment{CSLReferences}[2] % #1 hanging-indent, #2 entry-spacing
 {\begin{list}{}{%
  \setlength{\itemindent}{0pt}
  \setlength{\leftmargin}{0pt}
  \setlength{\parsep}{0pt}
  % turn on hanging indent if param 1 is 1
  \ifodd #1
   \setlength{\leftmargin}{\cslhangindent}
   \setlength{\itemindent}{-1\cslhangindent}
  \fi
  % set entry spacing
  \setlength{\itemsep}{#2\baselineskip}}}
 {\end{list}}
\usepackage{calc}
\newcommand{\CSLBlock}[1]{\hfill\break\parbox[t]{\linewidth}{\strut\ignorespaces#1\strut}}
\newcommand{\CSLLeftMargin}[1]{\parbox[t]{\csllabelwidth}{\strut#1\strut}}
\newcommand{\CSLRightInline}[1]{\parbox[t]{\linewidth - \csllabelwidth}{\strut#1\strut}}
\newcommand{\CSLIndent}[1]{\hspace{\cslhangindent}#1}

\KOMAoption{captions}{tableheading}
\makeatletter
\@ifpackageloaded{caption}{}{\usepackage{caption}}
\AtBeginDocument{%
\ifdefined\contentsname
  \renewcommand*\contentsname{Table of contents}
\else
  \newcommand\contentsname{Table of contents}
\fi
\ifdefined\listfigurename
  \renewcommand*\listfigurename{List of Figures}
\else
  \newcommand\listfigurename{List of Figures}
\fi
\ifdefined\listtablename
  \renewcommand*\listtablename{List of Tables}
\else
  \newcommand\listtablename{List of Tables}
\fi
\ifdefined\figurename
  \renewcommand*\figurename{Figure}
\else
  \newcommand\figurename{Figure}
\fi
\ifdefined\tablename
  \renewcommand*\tablename{Table}
\else
  \newcommand\tablename{Table}
\fi
}
\@ifpackageloaded{float}{}{\usepackage{float}}
\floatstyle{ruled}
\@ifundefined{c@chapter}{\newfloat{codelisting}{h}{lop}}{\newfloat{codelisting}{h}{lop}[chapter]}
\floatname{codelisting}{Listing}
\newcommand*\listoflistings{\listof{codelisting}{List of Listings}}
\makeatother
\makeatletter
\makeatother
\makeatletter
\@ifpackageloaded{caption}{}{\usepackage{caption}}
\@ifpackageloaded{subcaption}{}{\usepackage{subcaption}}
\makeatother

\ifLuaTeX
  \usepackage{selnolig}  % disable illegal ligatures
\fi
\usepackage{bookmark}

\IfFileExists{xurl.sty}{\usepackage{xurl}}{} % add URL line breaks if available
\urlstyle{same} % disable monospaced font for URLs
\hypersetup{
  pdftitle={Not Normal: A Simulation Study Comparing Effect Sizes for Skewed Psychological Data},
  pdfauthor={Ambra Perugini; Giulia Calignano; Massimiliano Pastore},
  pdfkeywords={Effect
sizes, Simulation, Overlapping, Assumptions, Non-normality},
  colorlinks=true,
  linkcolor={blue},
  filecolor={Maroon},
  citecolor={Blue},
  urlcolor={Blue},
  pdfcreator={LaTeX via pandoc}}


\title{Not Normal: A Simulation Study Comparing Effect Sizes for Skewed
Psychological Data}
\author{Ambra Perugini \and Giulia Calignano \and Massimiliano Pastore}
\date{}

\begin{document}
\maketitle
\begin{abstract}
Effect sizes are widely used in psychological science, but with a mere
focus on mean differences. This study compares three mean-based indices:
Cohen's \(d\), the Common Language Effect Size (CLES), the parametric
overlap coefficient derived from d (\(\eta_p\)), to the nonparametric
Overlapping Index (\(\eta\)), which focuses on the entire distributional
differences. Using skew-normal distributions, to generate controlled
scenarios of violations of normality and variance homogeneity, and
symmetry, we systematically manipulated mean differences, variance
ratios, skewness, and sample size to evaluate each index in terms of
Relative Mean Bias, Normalized Root Mean Square Error, and 90\%
Coverage. Cohen's \(d\), CLES, and \(\eta_p\) showed near-perfect
descriptive correlations, but their performance differed substantially.
Cohen's \(d\) consistently displayed low bias, high precision, and
accurate coverage across all scenarios, whereas CLES and \(\eta_p\)
exhibited substantial bias and low coverage, particularly under skewness
and heteroscedasticity. The nonparametric index \(\eta\) remained
unbiased under shape differences and variance heterogeneity but
performed less reliably when the populations truly overlapped. These
results demonstrate that effect-size indices derived from \(d\) are not
interchangeable, and that high empirical correlations do not guarantee
same precision. Cohen's \(d\) remains the most robust estimator of a
location difference, whereas \(\eta\) provides a more complete view of
the entire distribution. We argue that researchers should select
effect-sizes based on their statistical properties and propose a shift
towards interpreting effects in light of the full distribution, rather
than through mean-based conventions.
\end{abstract}


\section{Introduction}\label{introduction}

Researchers often behave as if a single number could capture the entire
contrast between two empirical worlds. In psychological research,
standardized effect size metrics are ubiquitous. As suggested by Lakens
(2013), Cohen's d is ``the most widely used effect size in psychological
research'' and is routinely used in meta-analyses aggregating group
differences. As proposed by McGraw \& Wong (1992) and Mastrich (2021),
the Common Language Effect Size (CLES) is ``an appealing index of effect
size,'' offering an intuitive interpretation in terms of probabilities
of superiority. In their standard parametric form, both Cohen's d and
CLES answer to the distance between means under assumptions of smooth,
symmetric, and homoscedastic distributions. Under those conditions,
variance does not rebel, and tails do not bend.

Yet the distributions underlying psychological data rarely obey this
ideal. Even in purely simulated scenarios, once asymmetry, variance
heterogeneity, and modest or large mean shifts are introduced, the data
no longer present a single clean contrast in location. Skewed
distributions, heavy right tails, and unequal dispersions create
conditions where the ``same'' effect, expressed as a standardized mean
difference, may emerge from very different shapes. Metrics that compress
the entire contrast into a standardized difference in means can then
misrepresent how groups truly diverge: they may be accurate about a
central shift while remaining blind to the structure of the tails or to
the degree of overlap between distributions.

As defined by Pastore \& Calcagnì (2019), the Overlapping Index
(\(\eta\)) formalizes a distribution-sensitive perspective by
quantifying how much two empirical density functions intersect. Unlike
parametric overlap measures derived from Cohen's \(d\) (such as the
Parametric Overlap, \(\eta_p\)), \(\eta\) does not rely on normality or
equal variances and can register differences that live in skewness,
kurtosis, or tail behavior. In this sense, Cohen's d, CLES, \(\eta_p\),
and \(\eta\) can be regarded as four different storytellers: some
summarizing only where the centers lie, others attending to how far and
how often the distributions diverge in their shapes.

The present work compares these four metrics that is, Cohen's \(d\),
CLES, \(\eta_p\), and \(\eta\) across a controlled set of non-normality
scenarios generated from skew-normal distributions. By systematically
varying mean differences, variance ratios, skewness, and sample sizes,
and by computing all four indices on the same simulated datasets, we
examine both how strongly they agree in practice and how well they
recover the underlying population contrasts. We evaluate their behavior
through Relative Mean Bias, Normalized Root Mean Square Error, and
Coverage, thereby assessing not only the magnitude each index reports
but also its bias, stability, and inferential reliability under
violations of normality and homoscedasticity. These results reveal when
classical, mean-based storytellers (\(d\), CLES, \(\eta_p\)) essentially
retell the same narrative, and when the distribution-based \(\eta\)
instead traces a distinct contour of difference, grounded in how much
the two distributions truly overlap rather than in how far apart their
means appear.

Of notice, the simulation-based approach is useful because it allows the
evaluation of effect-size indices under fully controlled conditions,
where the data-generating mechanisms are explicitly known and can be
manipulated independently. As emphasized in the broader methodological
literature on model assessment and inference (Akaike, 1973; Burnham \&
Anderson, 2002; Wagenmakers \& Farrell, 2004), the simulation provides a
principled way to examine how estimators behave when theoretical
assumptions are systematically violated. By varying mean differences,
variance ratios, skewness, kurtosis, and sample sizes, simulations make
it possible to quantify bias, precision, and robustness in ways that
real datasets that are necessarily, limited, idiosyncratic, and often
confounded cannot support. This framework exposes each metric to the
kinds of distributional structures that typify psychological data,
revealing where classical indices remain stable and where they fail. In
this sense, the simulation functions as an empirical stress test,
offering a clearer view of the operating characteristics of Cohen's
\(d\), CLES, \(\eta_p\), and \(\eta\) across the proposed landscape of
plausible data-generating scenarios.

\section{The Simulation Study}\label{the-simulation-study}

\subsection{Data Generation}\label{data-generation}

The simulation generates data sets from density distributions in
different scenarios. In particular, we used the Skew-Normal distribution
(Azzalini, 1985), which is defined in the following way: given
\(\xi \in \mathbb{R}\), \(\omega \in \mathbb{R}^{+}\) and
\(\alpha \in \mathbb{R}\), then for \(y \in \mathbb{R}\) we have

\begin{equation}\phantomsection\label{eq-skewnormal}{
\mathcal{SN}(y|\xi, \omega, \alpha)
  = \frac{1}{\omega \sqrt{2\pi}}
    \exp\left[
       -\frac{1}{2}\left( \frac{y-\xi}{\omega} \right)^2
    \right]
    \left[
      1 + \text{erf}\left(
           \alpha \left( \frac{y-\xi}{\omega\sqrt{2}} \right)
      \right)
    \right]
}\end{equation}

in which

\begin{equation}\phantomsection\label{eq-erf}{
\text{erf}(z)
  = \frac{2}{\sqrt{\pi}} \int_{0}^{z} e^{-t^2}\, dt
}\end{equation}

is the \emph{error function}.

\(\xi\) is the location parameter, \(\omega\) is the scale parameter and
\(\alpha\) is related to the skewness of the distribution. The case in
which \(\xi = 0\), \(\omega = 1\) and \(\alpha = 0\) corresponds to the
standard normal distribution. This distribution has been chosen to
generate the data as it allows to model both the distance between means
(the effect size), symmetry and variance. Consequently, it is optimal
for our purpose as it allows us to have control over the parameters of
mean, variance, skewness and kurtosis.

Mean and variance of the Skew-Normal are respectively:

\begin{equation}\phantomsection\label{eq-musigmaSN}{
\begin{aligned}
\mu &= \xi + \omega \gamma \sqrt{2/\pi} \\
\sigma^2 &= \omega^2 \left[1 - \frac{2\gamma^2}{\pi}\right]
\end{aligned}
}\end{equation}

in which \(\gamma = \alpha / \sqrt{1 + \alpha^2}\).

Based on Equation~\ref{eq-musigmaSN} we can determine the values to
assign to the parameters \(\xi\) and \(\omega\) as functions of \(\mu\)
and \(\sigma\) with:

\begin{equation}\phantomsection\label{eq-xiomegaSN}{
\begin{aligned}
\xi &= \mu - \omega \gamma \sqrt{2/\pi} \\
\omega &= \sqrt{\frac{\sigma^2}{1 - (2\gamma^2)/\pi}}
\end{aligned}
}\end{equation}

\subsection{Simulation design}\label{simulation-design}

The experiment consists of comparing two samples. We generated the first
sample from \(\mathcal{SN}(0,1,0)\), which corresponds to a Standard
Normal distribution with mean 0 and variance 1. We generated the second
sample from \(\mathcal{SN}(\xi,\omega,\alpha)\), where the parameters
are manipulated to produce specific differences in means (\(\delta\)),
standard deviations (\(\sigma\)), and skewness (\(\alpha\)) between the
two populations.

Overall, the experimental design comprises 8 scenarios. In each
scenario, the first sample remains fixed as
\(\mathcal{SN}(0,1,0) \equiv N(0,1)\), while the second sample is
defined by the following parameter settings:

\begin{itemize}
\item
  \(\mu\) = 0, 2; mean of the second population. As the first sample
  always comes from a population with a mean of 0, this parameter ---
  the mean of the population from which the second sample is drawn ---
  also represents the true difference between the means. Consequently,
  in the following sections of the paper we refers to this parameter as
  \(\delta\);
\item
  \(\sigma\) = 1, 5; standard deviation of the second population;
\item
  \(\alpha\) = 0, 10; asymmetry varies in the second group;
\item
  \(n\) = 10, 50, 100, 300, 500, 1000; sample size, equal in the two
  samples.
\end{itemize}

Figure~\ref{fig-simulationz} displays the experimental scenarios: the
black curves correspond to the first population, consistently modeled as
\(\mathcal{SN}(0,1,0) \equiv N(0,1)\), whereas the red curves represent
the second population \(\mathcal{SN}(\xi,\omega,\alpha)\). The vertical
dashed lines indicate the means of these second distributions and also
the mean difference, \(\delta\).

\begin{figure}[!t]

\centering{

\includegraphics{draft04_files/figure-latex/fig-simulationz-1.pdf}

}

\caption{\label{fig-simulationz}Representation of the simulation. In
each panel, the black curves correspond to the first population
\(\mathcal{SN}(0,1,0) \equiv N(0,1)\), while the red curves correspond
to the second population \(\mathcal{SN}(\xi,\omega,\alpha)\).}

\end{figure}%

For each of the 8 conditions we generated 2000 couples of independent
samples and performed the analysis. On these generated data, for each
combination \(\delta\times\sigma\times\alpha\times n\), we computed the
following effect sizes:

\begin{itemize}
\item
  Cohen's \(d\)
\item
  Common Language Effect Size (CLES).
\item
  Two kinds of Overlapping Coefficient:

  \begin{enumerate}
  \def\labelenumi{\arabic{enumi}.}
  \tightlist
  \item
    \emph{Parametric} (\(\eta_p\))
  \item
    \emph{Non-parametric} (\(\eta\)).
  \end{enumerate}
\end{itemize}

\subsection{Definition of Effect
sizes}\label{definition-of-effect-sizes}

\subsubsection{Cohen's d:}\label{cohens-d}

The Cohen's \(d\) is an effect size that standardizes differences in
means between groups or conditions (Cohen, 1988). This index
theoretically ranges from 0 to \(+\infty\) (for differences in one
direction) or \(-\infty\) to 0 (for the opposite direction), but in
practice, values rarely exceed a range of -3 to +3. A value of 0
indicates identical group means, whereas larger absolute values indicate
greater centroid distance. In the simulated scenarios, it is appropriate
to apply the between groups formula:

\begin{equation}\phantomsection\label{eq-d}{
d = \frac{\bar{X}_1 - \bar{X}_2}{s_p}
}\end{equation}

where \({s_p}\) is the pooled standard deviation defined as
\(s_p = \sqrt{\frac{(n_1 - 1)s_1^{2} + (n_2 - 1)s_2^{2}}{n_1 + n_2 - 2}}\).

\subsubsection{Common Language Effect Size
(CLES):}\label{common-language-effect-size-cles}

CLES represents the probability that a randomly picked individual from
the group with the higher mean will exceed a randomly picked individual
from the other group (McGraw \& Wong, 1992). This index ranges from 0 to
1 and expresses the probability that a randomly selected observation
from one group will exceed a randomly selected observation from another.
A value of 0.50 indicates complete overlap (no effect), values above
0.50 favor one group, and values below 0.50 favor the other, and
calculated from data the formula is:

\begin{equation}\phantomsection\label{eq-CLdata}{
CL = \frac{N_1}{N_1+N_0}
}\end{equation}

in which \(N_1\) indicates the number of ordered pairs in which the
first observation is greater than the second, and \(N_0\) the number of
ordered pairs in which the first observation is smaller than the second.

It can be also calculated as a function of Cohen's \(d\), \(\Phi\) is
the normal cumulative distribution function:

\begin{equation}\phantomsection\label{eq-CL}{
CL = \Phi(|d|/\sqrt{2})
}\end{equation}

\subsubsection{Overlapping Coefficient:}\label{overlapping-coefficient}

The concept of overlap refers to the shared area between two or more
probability density functions. It provides an intuitive measure for
assessing the degree of similarity or difference between samples or
populations characterized by their distributions. When two populations
(or samples) exhibit a high degree of overlap in their distribution
functions, it suggests a strong similarity between them. Conversely,
minimal overlap indicates substantial differences.

The \emph{parametric} version of the Overlapping coefficient can be
derived from Cohen's \(d\) as follows:

\begin{equation}\phantomsection\label{eq-OV}{
\eta_p = 2\Phi(-|d|/2)
}\end{equation} where \(\Phi\) is the normal cumulative distribution
function.

The \emph{non-parametric} version (\(\eta\)) is computed using a kernel
approach and serves as a distribution-free approximation of \(\eta_p\)
(for more details, see Pastore \& Calcagnì (2019)).

This coefficient varies from zero (indicating complete disjointness) to
one (indicating complete overlap) and can be considered an index of
similarity between groups or conditions.

\subsection{Performance Indexes for
comparison}\label{performance-indexes-for-comparison}

In order to compare the performance of the chosen indices, we need
indices that are not related to possible empirical values. We should
note that, whereas Cohen's \(d\) ranges over the real numbers, the other
three indices range from 0 to 1. Consequently, we used the following
performance indices:

\begin{itemize}
\item
  \emph{Relative mean bias} (RMB) captures the average proportional
  deviation between an estimator and its population value across Monte
  Carlo replications. It is defined as \[
  RMB = (1/B) \sum_b \left[ (\hat{\theta}_{bj} − \theta_j) / \theta_j \right].
  \] where \(\theta_j\) is the true parameter value in condition \(j\)
  and \(\hat{\theta}_{bj}\) is the estimated value in the \(b\)
  replications for condition \(j\). Because this ratio expresses error
  relative to the magnitude of the true parameter, it produces a
  scale-free measure of systematic distortion. Values close to zero
  indicate negligible bias, and values within the \(\pm 0.10\) interval
  are typically interpreted as acceptable. Whenever the true value
  \(\theta_j\) equals zero, a small constant \((\epsilon = 10^{-6})\) is
  substituted to prevent undefined divisions, ensuring numerical
  stability across replications.
\item
  \emph{Normalized root mean square error} (NRMSE) extends this
  scale-free approach to overall estimation error. It is simply derived
  from the Root Mean Square Error (RMSE) index dividing it by the
  empirical range of observed estimates, \[
  NRMSE = RMSE / (max(obs) − min(obs)),
  \] As with RMB, lower values indicate better performance.
\item
  \emph{Coverage} evaluates whether the inferential procedure is
  properly calibrated. It represents the proportion of simulated
  datasets whose 90\% confidence interval contains the true parameter
  value. High coverage indicates that the interval estimation procedure
  maintains its nominal reliability, independent of the scale of the
  underlying effect-size index.
\end{itemize}

Note that, other traditional indices such as Absolute Mean Bias, Mean
Square Error, and RMSE are less suitable in this context because they
remain scale-dependent. The relative and normalized metrics above
resolve this limitation by expressing accuracy on a common proportional
scale.

\section{Results}\label{results}

\subsection{Descriptive statistics}\label{descriptive-statistics}

In Figure~\ref{fig-correlations} reports the pairwise correlations among
the four effect size indices considered in the study (Cohen's \(d\),
CLES, \(\eta_p\), and \(\eta\)) across all 48 simulated conditions. The
diagonal panels display the univariate distribution of each index,
whereas the off-diagonal panels present the scatter plot (under the
diagonal) and pairwise correlation coefficients (above the diagonal);
different colors represent different sample sizes (n = 10, 50, 100, 300,
500, 1000).

\begin{figure}[!t]

\centering{

\includegraphics{draft04_files/figure-latex/fig-correlations-1.pdf}

}

\caption{\label{fig-correlations}Pairwise relationship between indexes.
The diagonal panels show the distribution of each index individually,
while the off-diagonal panels show the scatter plot (below the diagonal)
and the correlation coefficients (above the diagonal). Different colours
represent different sample sizes.}

\end{figure}%

The Figure~\ref{fig-correlations} have a very clear pattern, with
Cohen's \(d\), CLES and \(\eta_p\) showing very high correlations
between each other, consistently exceeding 0.95 with more than 50
observations and approaching 0.99 for larger sample sizes. In contrast,
\(\eta\) has only moderate correlation with those three indexes, ranging
between 0.40 and 0.50, suggesting that \(\eta\) reflects a more distinct
component of the data. Taken together, these findings indicate that
Cohen's \(d\), CLES and \(\eta_p\) capture highly overlapping
information about group differences, while \(\eta\) provides a partially
independent characterization of the effect, in line with the nature of
the effect size that takes into consideration the shape of the
distribution.

\subsection{Performance Indexes}\label{performance-indexes}

To evaluate the performance of the four effect sizes, we examined their
Relative Mean Bias (RMB), Normalized Root Mean Square Error (NRMSE), and
Coverage across all simulated scenarios.

\subsubsection{Relative Mean Bias (RMB)}\label{relative-mean-bias-rmb}

Figure~\ref{fig-RMB} shows the RMB for each index as a function of
sample size and under different experimental conditions. Overall,
\(\eta\) and Cohen's \(d\) display the lowest bias, with values
consistently close to zero across all conditions. CLES also exhibits
relatively low bias across scenarios with equal means, but shows biased
estimates when there is difference in means (\(\delta = 2\)). In the
opposite way, \(\eta_p\) shows unbiased estimates under same variance
(\(\sigma = 1\)) and biased estimates when variance is non-homogeneous
between the two groups (\(\sigma = 5\)).

\begin{figure}[!t]

\centering{

\includegraphics{draft04_files/figure-latex/fig-RMB-1.pdf}

}

\caption{\label{fig-RMB}Relative mean bias as a function of the
experimental conditions.}

\end{figure}%

\subsubsection{Normalized Root Mean Square
Error}\label{normalized-root-mean-square-error}

The NRMSE results (Figure~\ref{fig-NRMSE}) show that CLES displays much
larger NRMSE, dominating the scale of the plot and confirming its poorer
performance in terms of stability and precision when there are
differences in means between the two groups.

\begin{figure}[!t]

\centering{

\includegraphics{draft04_files/figure-latex/fig-NRMSE-1.pdf}

}

\caption{\label{fig-NRMSE}Normalized relative mean square error as a
function of the experimental conditions.}

\end{figure}%

For this reason, we also include a second plot (Figure~\ref{fig-NRMSEa})
in which CLES is excluded, as its performance is considerably worse than
that of the other indexes.

In this comparison, Cohen's \(d\) emerges as the most stable estimator,
with consistently lower NRMSE across all conditions. \(\eta_p\),
instead, shows high bias when homogeneity of variance is not respected,
but performs well under same variance conditions. \(\eta\) performs at
the level of Cohen's \(d\) in all condition but one, when the two groups
come from the same empirical distribution, reflecting the bounded nature
of the index.

\begin{figure}[!t]

\centering{

\includegraphics{draft04_files/figure-latex/fig-NRMSEa-1.pdf}

}

\caption{\label{fig-NRMSEa}Normalized relative mean square error
(without CLES) as a function of the experimental conditions.}

\end{figure}%

\subsubsection{Coverage}\label{coverage}

The 90\% Coverage results (Figure~\ref{fig-coverage}) show how Cohen's
\(d\) once more has good coverage in all scenarios, with \(\eta_p\) and
CLES having the worst coverage across scenarios and \(\eta\) performing
well in all scenarios, except in the two with \(\delta = 0\) and
\(\sigma = 1\).

\begin{figure}[!t]

\centering{

\includegraphics{draft04_files/figure-latex/fig-coverage-1.pdf}

}

\caption{\label{fig-coverage}Coverage 90\% as a function of the
experimental conditions.}

\end{figure}%

\section{Discussion}\label{discussion}

The simulation results demonstrate that the four effect-size indices
that is, Cohen's d, the Common Language Effect Size (CLES), the
Parametric Overlap derived from d (\(\eta_p\)), and the
distribution-free Overlapping Index \(\eta\) do not constitute
interchangeable descriptions of group or condition differences when
distributions deviate from normality. Across all manipulated conditions
(mean differences, variance ratios, skewness, and sample sizes), the
indices separate into two functional families: mean-based location
estimators (\(d\), CLES, \(\eta_p\)) and a distribution-sensitive
overlap estimator (\(\eta\)). This separation reflects that the indices
derived from \(d\) inherit all of its parametric assumptions, as
suggested by Peng \& Chen (2014). However, the performance metrics show
that this conceptual grouping does not imply equivalent inferential
properties: although the three \(d\)-based indices behave similarly
descriptively, they differ substantially in bias, stability, and
coverage.

Across the full simulation space, near-perfect correlations among \(d\),
CLES, and \(\eta_p\) (\(\rho\) \textgreater{} .93) confirm their shared
theoretical base. Each assumes that the primary quantity of interest is
a standardized shift in location, with dispersion treated as nuisance or
absorbed through pooled variance. CLES is mathematically linked to \(d\)
under symmetry, and \(\eta_p\) is a deterministic transformation of
\(d\) under normality, as corrected by Grice \& Barrett (2014). Thus,
under approximately symmetric or mildly skewed distributions, all three
indices encode the same latent structure, a rank-ordered monotonic
recoding of mean distance. Nevertheless, thier performances differ:
Cohen's \(d\) shows low bias across all conditions and high coverage,
indicating that it is able to identify the true parameter, but CLES and
\(\eta_p\) show biased estimates and low coverage in some conditions,
demonstrating that a high empirical correlation with \(d\) does not
guarantee accuracy or inferential validity.

Out of the group, \(\eta\) diverges systematically from classical
indices (\(r \in [ .42, .47 ]\)), not focusing structly on means but
capturing a broader distributional difference. Performing at the level
of Cohen's \(d\) for RMB, its weakest performance or RMSR emerges
exactly under strict normality, probably due to the nature of the index.
The index varies from \(0\), complete disjointment, and \(1\), complete
overlap, but two empirical populations will never overlap exactly even
with high sample size. Again, when the two populations parameters of
mean and variance are identical (\(\delta\) = 0, \(\sigma\) = 1),
particularly when they are completely the same, \(\eta\) shows low
coverage, reflecting sensitivity to kernel estimation when true overlap
is maximal.

In practical use, \(\eta\) relies on kernel-based density estimation,
which smooths fluctuations and mitigates outliers (Chen et al., 2021).
Importantly, the overlapping package permits adjustment of the kernel
function, introducing a controllable dimension of flexibility, what
might be called a tailored robustness check, that allows \(\eta\) to
adapt to diverse data structures.

At last, Cohen's \(d\) is confirmed to be a robust tool to asses
groups/conditions differences based on means. However, \(\eta\), rather
than standardized mean differences, quantifies shared mass between
empirical densities, thereby capturing distributional aspects ignored by
\(d\)-based measures. Simulation evidence confirms that \(\eta\) remains
unbiased even under strong variance heterogeneity (ratios of 3:1 or
4:1), unlike \(\eta_p\) and CLES, whose estimates become distorted under
such violations (Skidmore \& Thompson, 2013).

These findings are especially relevant in Cognitive Psychology and
Neuropsychology, where reaction-time (RT) distributions are typically
right-skewed and long-tailed. In such datasets, \(d\) may be a
meaningful effect size even when substantial portions of distributions
do not overlap. Still, it tells only part of the story, overlooking
shape characteristics preserved by \(\eta\) characteristics that can
also bear interpretable functional significance in cognitive and
neuropsychological frameworks. This empirical structure has been
repeatedly observed in memory and attention studies, as illustrated by
Osmon (2018). Under these conditions, critically using both \(d\) and
\(\eta\) provides a more realistic quantification of divergence, where
\(\eta\) is able to quantify differences that would be overlooked in
cases of equal means, isolating distributional differences that
mean-based indices obscure.

A similar principle applies in clinical psychology, where symptom
distributions are rarely normal (Rubio-Aparicio et al., 2018), and both
floor effects in non-clinical populations and heavy right tails in
clinical ones are common characteristics of symptom data (Haqiqatkhah et
al., 2023). These asymmetries inflate pooled variance, leading
traditional effect size measures such as Cohen's \(d\) and Common
Language Effect Size (CLES) to underestimate the true separations
between groups when distributions deviate from normality (Sun \& Cheung,
2020). In such cases, Cohen's \(d\) would still be able to recover the
true mean difference, but at the cost of theoretical vulnerability, not
correctly identifying differences in skewness and lack of variance
heterogeneity, which are valuable aspects of the data structure. In such
cases their combination provides a balanced framework for inference
(Hanel \& Mehler, 2019).

Here, \(\eta\) enters not as a rival, but as a companion. Where \(d\)
fixes its gaze upon the center, \(\eta\) listens to the contour to the
way distributions overlap, diverge, and rejoin across the span of shape
variation (Hedges \& Olkin, 2016). In doing so, it restores something
that \(d\) leaves silent: the shape of difference. Their relationship is
not one of substitution but of complementarity as two instruments,
distinct in tone yet tuned to the same melody.

Thus, the lesson is not that one index should prevail, but that both
deserve to speak. \(d\) translates difference into the language of
averages; \(\eta\) translates it into that of a broader concept of
similarity between distributions. And somewhere between those two
dialects between what is measured and what is experienced lies the truer
inference, the quiet understanding that every distribution in
psychology, like every life, bends a little away from symmetry.

\section{Conclusion}\label{conclusion}

Effect sizes should not be treated as interchangeable. When the
scientific question concerns a simple shift in location and assumptions
are reasonably satisfied, \(d\) represents an adequate and robust index,
with very little bias and high coverage of the true parameter. But when
such assumptions are not respected, and the question concerns
distributional separation like tail behavior, heterogeneity,
multimodality, or skew-driven divergence \(\eta\) provides the useful
and complementary description. An index as \(\eta\) can capture
differences that would otherwise be overlooked when merely focusing on
means. With the support of an adequate data visualization process, those
differences become even more clear to the researcher and allow for a
more tailored interpretation.

The simulation approach comparing effect size indexes employed here
follows established methodological frameworks advocating controlled
assumption violations to evaluate estimator performance (Zhao et al.,
2020). The core emphasis of the present study here, is that no single
effect-size metric suffices; robust interpretation depends on viewing
several indices together and anchoring them to the shape and structure
of the observed data. Simulations offer what empirical data cannot: a
controlled environment for systematic stress testing. At the same time,
it reflects the very distortions that shape real-world evidence,
allowing us to understand how psychological effect-size indices hold or
fail under those pressures.

\section{References}\label{references}

\phantomsection\label{refs}
\begin{CSLReferences}{1}{0}
\bibitem[\citeproctext]{ref-akaike:1973}
Akaike, H. (1973). {Information theory and an extension of the maximum
likelihood principle}. In B. N. Petrov \& F. Csaki (Eds.), \emph{{Second
International Symposium on Information Theory}} (pp. 267--281).

\bibitem[\citeproctext]{ref-azzalini:1985}
Azzalini, A. (1985). A class of distributions which includes the normal
ones. \emph{Scandinavian Journal of Statistics}, 171--178.

\bibitem[\citeproctext]{ref-burnham+anderson:2002}
Burnham, K. P., \& Anderson, D. R. (2002). \emph{{Model {S}election and
{M}ultimodel {I}nference: {A} {P}ractical {I}nformation-{T}heoretic
{A}pproach}} (2nd ed.). Springer Science \& Business Media.

\bibitem[\citeproctext]{ref-Chen:2021}
Chen, X., Trafimow, D., Wang, T., Tong, T., \& Wang, C. (2021). {The APP
Procedure for Estimating the Cohen's Effect Size}. \emph{Asian Journal
of Economics and Banking}.

\bibitem[\citeproctext]{ref-Cohen:1988}
Cohen, J. (1988). \emph{{Statistical power analysis for the behavioral
sciences}} (2nd ed.). Lawrence Erlbaum Associates, Hillsdale, NJ.

\bibitem[\citeproctext]{ref-Grice:2014}
Grice, J., \& Barrett, P. (2014). {A Note on Cohen's Overlapping
Proportions of Normal Distributions}. \emph{Psychological Reports},
\emph{115}, 741--747.

\bibitem[\citeproctext]{ref-Hanel:2019}
Hanel, P., \& Mehler, D. (2019). {Beyond Reporting Statistical
Significance: Identifying Informative Effect Sizes to Improve Scientific
Communication}. \emph{Public Understanding of Science}, \emph{28},
468--485.

\bibitem[\citeproctext]{ref-Haqiqatkhah:2023}
Haqiqatkhah, M. M., Ryan, O., \& Hamaker, E. (2023). {Skewness and
Staging: Does the Floor Effect Induce Bias in Multilevel AR(1) Models?}
\emph{Multivariate Behavioral Research}, \emph{59}, 289--319.
\url{https://doi.org/10.1080/00273171.2023.2254769}

\bibitem[\citeproctext]{ref-Hedges:2016}
Hedges, L., \& Olkin, I. (2016). {Overlap between treatment and control
distributions as an effect size measure in experiments}.
\emph{Psychological Methods}, \emph{21}(1), 61--68.
\url{https://doi.org/10.1037/met0000042}

\bibitem[\citeproctext]{ref-lakens:2013}
Lakens, D. (2013). {Calculating and reporting effect sizes to facilitate
cumulative science: A practical primer for t-tests and ANOVAs}.
\emph{Frontiers in Psychology}, \emph{4}, 863.
\url{https://doi.org/10.3389/fpsyg.2013.00863}

\bibitem[\citeproctext]{ref-mastrich:2021}
Mastrich, Z. (2021). {A review of common language effect size indicators
to facilitate research interpretability}. \emph{Frontiers in
Psychology}. \url{https://doi.org/10.3389/fpsyg.2021.648770}

\bibitem[\citeproctext]{ref-mcgraw+wong:1992}
McGraw, K. O., \& Wong, S. P. (1992). {A common language effect size
statistic}. \emph{Psychological Bulletin}, \emph{111}(2), 361--365.
\url{https://doi.org/10.1037/0033-2909.111.2.361}

\bibitem[\citeproctext]{ref-Osmon:2018}
Osmon, K., D. C. (2018). {Non-Gaussian distributional analyses of
reaction times (RT): Improvements that increase efficacy of RT tasks for
describing cognitive processes.} \emph{Neuropsychology Review},
\emph{28}, 359--376.

\bibitem[\citeproctext]{ref-pastore+calcagni:2019}
Pastore, M., \& Calcagnì, A. (2019). {Measuring distribution
similarities between samples: A distribution-free overlapping index}.
\emph{Frontiers in Psychology}, \emph{10}, 1089.
\url{https://doi.org/10.3389/fpsyg.2019.01089}

\bibitem[\citeproctext]{ref-Peng:2014}
Peng, C. Y., \& Chen, L. T. (2014). {Beyond Cohen's d: Alternative
Effect Size Measures for Between-Subject Designs}. \emph{The Journal of
Experimental Education}, \emph{82}, 22--50.

\bibitem[\citeproctext]{ref-RubioAparicio:2018}
Rubio-Aparicio, M., Marín-Martínez, F., Sánchez-Meca, J., \&
López-López, J. (2018). {A methodological review of meta-analyses of the
effectiveness of clinical psychology treatments}. \emph{Behavior
Research Methods}, \emph{50}, 2057--2073.
\url{https://doi.org/10.3758/s13428-017-0973-8}

\bibitem[\citeproctext]{ref-Skidmore:2013}
Skidmore, S. T., \& Thompson, B. (2013). {Bias and Precision of Some
Classical ANOVA Effect Sizes When Assumptions Are Violated}.
\emph{Behavior Research Methods}, \emph{45}, 536--546.

\bibitem[\citeproctext]{ref-Sun:2020}
Sun, R., \& Cheung, S. (2020). {The influence of nonnormality from
primary studies on the standardized mean difference in meta-analysis}.
\emph{Behavior Research Methods}, 1--16.
\url{https://doi.org/10.3758/s13428-019-01334-x}

\bibitem[\citeproctext]{ref-wagenmakers+farrell:2004}
Wagenmakers, E.-J., \& Farrell, S. (2004). {AIC} model selection using
{A}kaike weights. \emph{Psychonomic Bulletin \& Review}, \emph{11}(1),
192--196.

\bibitem[\citeproctext]{ref-Zhao:2020}
Zhao, J., Sim, S., \& Kim, H.-M. (2020). {Statistical Inference of Some
Effect Sizes}. \emph{Journal of the Korean Statistical Society},
\emph{49}, 976--1007.

\end{CSLReferences}




\end{document}
